% Discussion and limitations. Numbers are macros from generated/numbers.tex.

\paragraph{What the results support.} On synthetic clouds drawn from the 13-family generative model, a network trained only on synthetic data identifies the family and recovers the shape better than a classical pipeline that fits every family and selects with geometric AIC (H1, H2), and its probabilities are well calibrated after temperature scaling. This is a \emph{systematic} claim within the synthetic distribution (\nBone{} paired clouds in 156 generator chunks for H1 and H2, 20 per family and cell over all point-count and noise bins), with paired cluster-bootstrap CIs. It is a claim about the matched synthetic prior, which the classifier has learned and the classical pipeline does not use. The preregistered B1 overstates the gap: geometric AIC as implemented can never select the blob. The comparison therefore also holds on non-blob clouds and against the post-hoc repair B1-fixed (Sec.~\ref{sec:results}). The ablation A5 attributes a large share of the advantage to model selection: given the same refined candidates, geometric AIC selects worse than the learned classifier, also on non-blob clouds.

\paragraph{What the results do not support.} The fitter is not yet a reliable replacement for by-eye shape calls on real activations. Coarse agreement with the literature (H3) is not above a permutation null. Post-hoc inspection suggests causes for the coarse misses: a curve outside the vocabulary (GPT-2 positions), a top-3 projection of higher-dimensional structure (years), clouds unlike any training family (age, true/false statements) and an unstable label on an 8-point cloud. At the fine level, more than six discrete groups (weekdays, months) are threaded with a helix. These explanations were not tested. In each case the classifier still returns a label, because a closed-set classifier must pick one of its 13 families; its calibrated probability can be low (\pRsixllamayears{} for years) or high (\pRthreetmsneight{} for the 8-point toy model). Smooth ordinal controls are frequently labelled helix (\hfourGpHelix{} of the Gaussian-process controls, post hoc), so a helix label from the classifier alone is not evidence of a helix. The preregistered two-part rule (probability ratio and held-out residual) is necessary but not sufficient: it rejected the helix claim for the one Gaussian-process control on which it was run (the other draws got only the full-cloud fit) but accepted it for \helixRuleAccepted{}, which the literature describes as circles.

\paragraph{Implications for practice.} (1) Report the shape of the top-3 projection together with the explained variance and the $\lambda_3/\lambda_4$ gap; a low gap means the 3D picture is a slice. (2) Treat ``helix'' as a hypothesis to be tested with held-out data against the PCA ringing artifact \citep{novembre2008horseshoe,antognini2018walks,shinn2023phantom} and against discrete clusters or a circle. (3) Controls with random weights are not structureless when inputs differ in token identity; a token-identity-matched control is the right null for prompt-based datasets. (4) Exclude sequence position 0 (BOS or the first token) when averaging activations: the starting token carries massive activations and attracts attention as a sink \citep{sun2024massive,xiao2024sinks}, which can dominate class means.

\paragraph{Limitations.}
\label{sec:limitations}
\begin{itemize}
\item \textbf{One seed, no trained ablations yet.} All results come from one trained model (\texttt{main\_s0}); the seed, learning-rate and inference-hyperparameter sweeps and the trained ablations A1, A2, A4, A6 (including H5a) are implemented but postponed to later stages (\texttt{reports/deviations.md}). Synthetic CIs therefore reflect test-sample variation only, not training variation.
\item \textbf{Baselines on a subset.} B1 and B2 were run on 20 clouds per family and cell (\nBone{} clouds) because B1 runs up to 61 refinements per cloud; all comparisons with baselines use this subset (H5b is paired on the full test set). B1 is also handicapped in three ways found after the run. (i)~Geometric AIC with the noise estimated from the best surface fit can never select the blob. (ii)~Its geometric-AIC residual $J$ is measured to the nearest of a finite set of shape samples (a $40\times40$ grid for surfaces), which adds discretisation error comparable to the lowest noise levels. (iii)~In batches containing 5- or 6-point clouds it skips 5- and 6-cluster candidates (\texttt{reports/deviations.md}). B1-fixed repairs all three post hoc; its blob threshold is one more fitted quantity, tuned on validation clouds only. The real-data B1 column still uses the preregistered B1.
\item \textbf{Dependent test clouds.} Clouds generated together share $N$ and the embedding draw. CIs therefore resample generator chunks; before the 2026-11-14 audit they resampled clouds and were too narrow. Because chunks are resampled across the design's family $\times$ point-count $\times$ noise cells, the reported CIs are conservative (Sec.~\ref{sec:setup}). The ``unambiguous'' subset is defined by B1's own fits and favours B1.
\item \textbf{Chamfer grid floor.} Both shapes enter the Chamfer distance through fixed evaluation samples (256 per curve, a $24\times24$ grid per surface). Against a 16 times denser sampling of the same shape, an evaluation sampling scores at most \floorCurveMax{} for curves but \floorSurfMin{}--\floorSurfMax{} for surfaces. Two equivalent parameterisations of a \floorSymFams{} (rotated about its axis) score up to \floorSymMax{}. Absolute Chamfer values of surfaces therefore partly measure the grid. The floor enters every system's score, and the rotation about a symmetry axis cannot be inferred from the points, so we expect it to add noise rather than favour one system (post hoc, audit 2026-11-22).
\item \textbf{Closed vocabulary.} 13 families with label margins between nested families; shapes outside the vocabulary are forced into one of them. There is no abstention beyond the OOD and gap flags.
\item \textbf{Synthetic prior.} Training data follow hand-designed families and nuisances; real activation clouds differ, for example in token-level clusters, and we did not test which difference causes the synthetic-to-real gap.
\item \textbf{Top-3 PCA only.} All statements concern the top-3 projection, as the task specifies; intrinsic manifolds can differ, and causal roles of these representations are not tested.
\item \textbf{Literature labels as ground truth.} H3 compares against published descriptions that were themselves made by eye, often with different models, layers or prompts (we use Llama-3.1-8B for all preregistered LLM datasets). Some are labels from a dataset table that differ from the source's own PCA view: the ``line'' for temperature and age and the ``helix'' for years of \citet{bhalla2026sae}. The grokked model was trained for fewer epochs than in \citet{nanda2023progress}, and one weekday prompt differs by one word from \citet{engels2025notall}.
\end{itemize}
